\documentclass{ximera}

\title{Introductie}
\author{Mila Vervoort}
\license{CC: 0}         % replace with an appropriate license, or set it in xmPreamble

\begin{document}
\begin{abstract}
    Samenvatting van de theorie over reactiekinetiek.
\end{abstract}
\maketitle
\label{xim:Introductie analysemethode}

\section*{Analyse en verwerking van kinetische metingen}

\subsection*{Introductie}

Eigenschappen van chemische reacties zijn vaak ongekend en niet eenvoudig te verkrijgen.

Soms, voornamelijk voor elementaire gasfase reacties, zijn er benaderingen mogelijk
of kan men quantumchemische berekeningen uitvoeren.

De meest gebruikte methode om chemische kinetica te bepalen is echter via
\textbf{experimentele metingen}. In dit hoofdstuk behandelen we de analyse
en verwerking van experimentele metingen om kinetische parameters te bepalen.

We gaan uit van een functionele vorm voor de reactiesnelheid:
\[-R_A = f(C_A)g(T)
\]

\subsection*{Kinetische metingen in een batch reactor}

Deze methode wordt vooral gebruikt voor homogene reacties. We veronderstellen een
isotherme reactor met constant volume.

Om een uitdrukking voor de reactiesnelheid te bepalen, worden concentratiemetingen op verschillende tijdstippen uitgevoerd.

Soms wordt de concentratie gekoppeld aan andere meetbare eigenschappen, zoals:

\begin{itemize}
\item thermische conductiviteit
\item viscositeit
\item spectrometrie
\end{itemize}

\begin{example}
Voor gasfase reacties kan men concentraties eenvoudig bepalen door de totale druk en de partieeldrukken te meten.

\[
\begin{aligned}
aA + bB &\rightarrow rR + sS \\[6pt]
N_A &= N_{A0} - ax \\
N &= N_0 + (r + s - a - b)x = N_0 + \Delta n\,x \\[6pt]
C_A &= \frac{N_A}{V} = \frac{p_A}{RT} \\
p_A &= p_{A0} - \frac{a}{\Delta n}(p - p_0)
\end{aligned}
\]
\end{example}

Wanneer kinetische metingen de concentratie $C_A$ op verschillende tijdstippen $t$ opleveren, 
kan de reactiekinetiek hieruit worden afgeleid. 

Hiervoor kunnen verschillende methoden worden gebruikt:
\begin{itemize}
    \item Integrale analysemethode
    \item Differentiële analysemethode
\end{itemize}


Anders:
\subsection*{Inleiding}

Om een chemische reactor te kunnen ontwerpen, moet men de \textbf{reactiesnelheid}
van de betrokken reactie kennen. Deze wordt beschreven door een
\textbf{snelheidsvergelijking}, die aangeeft hoe snel een reactie verloopt als functie
van bijvoorbeeld concentratie en temperatuur.

De functionele vorm van deze vergelijking kan soms afgeleid worden uit
theoretische modellen, maar wordt vaak bepaald via \textbf{empirische fitting}
van experimentele gegevens. De numerieke waarden van de kinetische parameters
(bijvoorbeeld snelheidsconstanten) worden vrijwel altijd bepaald uit
\textbf{experimentele metingen}.

Kinetische experimenten hebben meestal als doel:
\begin{itemize}
\item de afhankelijkheid van de reactiesnelheid van de \textbf{concentraties} te bepalen
\item de \textbf{temperatuurafhankelijkheid} van de reactiesnelheid te onderzoeken
\end{itemize}

Door deze informatie te combineren kan men de volledige
\textbf{reactiesnelheidsvergelijking} opstellen.

\subsection*{Reactoren voor kinetische metingen}

Voor het uitvoeren van kinetische experimenten worden vooral twee types reactoren gebruikt:

\begin{itemize}
\item \textbf{Batchreactoren}: vaak gebruikt voor homogene reacties. De samenstelling
van het mengsel verandert in de tijd terwijl er geen instroom of uitstroom is.
\item \textbf{Continue reactoren}: bijvoorbeeld kuipreactoren (CSTR) of buisreactoren
(PFR), die vaak gebruikt worden bij heterogene of katalytische systemen.
\end{itemize}

De verkregen experimentele data kan vervolgens geanalyseerd worden met verschillende
methoden, zoals de \textbf{integrale methode} en de \textbf{differentiële methode},
om de kinetische parameters van de reactie te bepalen.


\end{document}